\section{Oriented reciprocal pooling at different bases}

The next reduction is independent of the center construction. It
starts from package 173's unchanged signed Pell source, whose frame
already gives $|A|>W^W$ and bounds the squares of every actual base,
requested output, and nominal index by $W$ \cite{source173}.
The Pell polynomials are evaluated at signed integer $A$ using their
polynomial recurrences. For signed indices use $X_{-k}(A)=X_k(A)$ and
$Y_{-k}(A)=-Y_k(A)$. These are the coefficients of the corresponding
positive or negative powers of the unit $A+s$, $s^2=A^2-1$.

\begin{theorem}[Oriented pooling with inactive zero bases]\label{thm:pool}
Let $W\in\Z$, $W\ge258$, let $A\in\Z$ satisfy $|A|>W^W$, and let at most five
integer bases $n_j$ be given. There are at most 64 requests in total.
Each request in block $j$ consists of integers $(m_i,k_i)$ and canonical
signed Pell coordinates $(X_i,Y_i)$ at parameter $A$ and index $k_i$.
Assume
\[
 n_j^2<W,\qquad m_i^2<W,\qquad k_i^2<W.
\]
Define
\begin{align*}
 T_j&=2An_j-n_j^2-1,&D_i&=X_i-AY_i,\\
 S_j&=\sum_{i\in j}\bigl((m_i-D_i)^2-Y_i^2\bigr),&
 H_j&=\sum_{i\in j}Y_i(m_i-X_i),\\
 E_j&=S_j-2n_jH_j.
\end{align*}
All $T_j$ are nonzero. Then
\begin{equation}\label{eq:orientedintegrality}
 \sum_j\frac{n_jE_j}{T_j}\in\Z
\end{equation}
if and only if every request in every nonzero-base block satisfies
$m_i=n_j^{k_i}$, initially interpreted in $\Q$. For $|n_j|\ge2$ this
also forces $k_i\ge0$. No conclusion is asserted about a zero-base
block's outputs; its term in \eqref{eq:orientedintegrality} is zero.
\end{theorem}

\begin{proof}
For $n_j=0$, $T_j=-1$. For $n_j\ne0$,
\[
 |T_j|\ge2|A|-W>W^W,
 \qquad\operatorname{sgn}\frac{n_j}{T_j}=\operatorname{sgn}A.
\]
Let $\delta_i=m_i-X_i+(A-n_j)Y_i$. Direct expansion gives
\begin{equation}\label{eq:erroridentity}
 E_j=\sum_{i\in j}\delta_i^2+T_j\sum_{i\in j}Y_i^2.
\end{equation}
Thus \eqref{eq:orientedintegrality} is equivalent to integrality of
$\sum_j n_j\sum_i\delta_i^2/T_j$.

We record why signed indices cause no gap. In the quadratic integer
ring with $s^2=A^2-1$, send $s$ to $n_j-A$ modulo $T_j$. This is a
ring homomorphism because
\[
 (n_j-A)^2-(A^2-1)=n_j^2-2An_j+1=-T_j.
\]
The unit $A+s$ maps to $n_j$, whose inverse is $2A-n_j$, since
$n_j(2A-n_j)=1\pmod{T_j}$. Therefore for every signed integer $k_i$,
\begin{equation}\label{eq:signedcongruence}
 X_i+(n_j-A)Y_i\equiv n_j^{k_i}\pmod{T_j}.
\end{equation}
Here negative powers on the right mean powers of this modular unit.

Put $b=\lfloor\sqrt{W-1}\rfloor$, so $b\ge16$. Let $J$ be the set
of nonzero-base blocks, $s=|J|\le5$, and put
\[
 C=\prod_{j\in J}n_j^b\ne0.
\]
The integer $C$ appears only in the proof. For a request in an active
block define
\[
 d_i=Cm_i-\frac{C}{n_j^b}n_j^{b+k_i}\in\Z.
\]
The exponents $b,b+k_i$ are nonnegative. Equation
\eqref{eq:signedcongruence} implies
$C\delta_i\equiv d_i\pmod{T_j}$, and hence
$C^2\delta_i^2-d_i^2$ is a multiple of $T_j$.
Multiplying the integral sum by $C^2$ now shows that
\[
 \Lambda=\sum_{j\in J}\frac{n_j\sum_{i\in j}d_i^2}{T_j}
\]
is an integer. This step clears each negative-index denominator before
any estimate is used.

Every relevant absolute value is at most $b$, so
\[
 |d_i|\le b^{sb+1}+b^{(s+1)b}\le2b^{(s+1)b}.
\]
Writing $r\le64$ for the total number of requests, we obtain
\begin{align*}
 \sum_{j\in J}|n_j|\sum_{i\in j}d_i^2
 &\le4r\,b^{2(s+1)b+1}\\
 &\le4r\,b^{12b+1}<W^{6b+2}\le W^W.
\end{align*}
The strict inequality follows from $4r\le256<W$ and $b^2<W$.
The last inequality uses $6b+2\le b^2+1\le W$, valid for $b\ge16$.
Since every denominator has absolute value greater than $W^W$,
$|\Lambda|<1$. All its nonzero summands have sign $\operatorname{sgn}A$.
Thus the integer $\Lambda$ vanishes, and every $d_i=0$. Dividing by
$C$ gives $m_i=n_j^{k_i}$ in $\Q$. If $|n_j|\ge2$ and $k_i<0$,
the integer equality $m_i n_j^{-k_i}=1$ is impossible.

Conversely, exact powers imply $T_j\mid\delta_i$ in every active
block. Equation~\eqref{eq:erroridentity} makes $E_j/T_j$ integral
there. Every inactive block contributes zero, proving
\eqref{eq:orientedintegrality}.
\end{proof}

\subsection{Clearing the denominators without adding a primitive}

Introduce one unrestricted integer $z$ and define the ordinary polynomial
\begin{equation}\label{eq:fullpool}
 \mathcal R=
 \sum_{j=0}^{4}n_jE_j\!\prod_{\substack{0\le h\le4\\h\ne j}}T_h
 -z\prod_{h=0}^{4}T_h.
\end{equation}
Since every $T_h\ne0$ under the retained frame, $\mathcal R=0$ is
equivalent to \eqref{eq:orientedintegrality}. The denominator notation
has disappeared from the actual equation. Existing multiplication wires
$Z_i=AY_i$ allow the completely literal definition
\[
 E_j=\sum_{i\in j}
 \left((m_i-D_i)^2-Y_i^2
       -2n_jY_i(m_i-D_i-Z_i)\right).
\]
Thus $E_j$ has degree at most three and $T_j$ degree two. Formula
\eqref{eq:fullpool} has degree at most twelve. Its emitted coefficient
list has a nonzero degree-12 part, so its square has exact degree 24.

For the two existing fixed-base pools, the old private variables
$K_j,L_j$ satisfy
\[
 K_j-2H_j-(2A-n_j)L_j=0,\qquad S_j-n_jK_j+L_j=0.
\]
Adding $n_j$ times the first row to the second gives $E_j-T_jL_j$.
The construction verifies this identity directly against the literal
parent rows; the block membership is not inferred from descriptive
variable names.

\subsection{Activating every block without assuming a positive radix}

The five actual bases are
\[
 t,\quad B,\quad BH,\quad q,\quad q+1,
\]
with respectively $1,1,1,3,10$ requests. All their canonical signed
Pell coordinates and the frame bounds are recovered before any deleted
power row is used. The unchanged source supplies
\begin{gather*}
 t=128+4h^2,\qquad K=t^4,\qquad F\text{ requested as }t^K,\\
 C=2c,\qquad B\mu=C+F+1,\qquad
 H=2+4\sum(\text{actual CRT representatives}),\\
 q\text{ requested as }(BH)^x.
\end{gather*}
The following order is essential.
\begin{enumerate}
\item Since $t\ge128$, its block is active. Theorem~\ref{thm:pool}
gives $F=t^K$. In particular $F$ is even.
\item The literal $C=2c$ is even. Hence $C+F+1$ is odd and nonzero.
The retained equality $B\mu=C+F+1$ forces $B$ odd and nonzero, without
using the sign of $C$.
\item $H\equiv2\pmod4$, so $H\ne0$ and $|H|\ge2$. The active
$BH$ block gives $x\ge0$ and $q=(BH)^x$.
\item If $x=0$, then $q=1$. If $x\ge1$, then $q$ is nonzero and
even. Thus $q\notin\{0,-1\}$, and both remaining bases are nonzero.
\end{enumerate}
All five blocks are now active and all original power assertions hold.
The proof has not changed a CRT representative, parity condition,
radix constructor, or shared residual to obtain this conclusion.

For a new zero, restore each singleton modulus as $T_j$ and its old
quotient as $-\delta_j/T_j$. For a fixed-base block restore
\[
 L_j=E_j/T_j,\qquad K_j=2H_j+(2A-n_j)L_j.
\]
These are integers by the recovered powers. Every retained coordinate
is unchanged. Conversely, at an old zero the old singleton quotient
$Q_j$ gives $\delta_j=-T_jQ_j$, so
$E_j/T_j=T_jQ_j^2+Y_j^2$. The explicit polynomial choice
\[
 z=\sum_{j=0}^{2}n_j(T_jQ_j^2+Y_j^2)
       +qL_q+(q+1)L_{q+1}
\]
satisfies the new certificate. Thus the transformation is pointwise
existential equivalence over the shared source coordinates.

\subsection{A degree-16 companion}

Use oriented pooling only on the first three bases, with quotient $z_0$:
\[
 \mathcal R_0=
 \sum_{j=0}^{2}n_jE_j\!\prod_{\substack{0\le h\le2\\h\ne j}}T_h
 -z_0\prod_{h=0}^{2}T_h.
\]
The preceding parity argument already makes $q,q+1$ nonzero. Two
nonzero consecutive integers have the same sign. Their denominators
therefore have the same sign, and the unweighted two-base certificate
from package 168 applies \cite{source168}:
\begin{equation}\label{eq:splitpool}
 \mathcal R_1=E_3T_4+E_4T_3-z_1T_3T_4.
\end{equation}
For completeness from an old zero take
$z_0=\sum_{j=0}^{2}n_j(T_jQ_j^2+Y_j^2)$ and
$z_1=L_q+L_{q+1}$. For soundness the same bounded-error proof uses
$C=(q(q+1))^b$ and the estimate
\[
 4\cdot13\,b^{6b}<W^{3b+1}\le W^W.
\]
There are no opposite-sign errors to cancel. The first new residual has
exact degree eight, the second exact degree five. Their squares and the
144 retained quadratic squares give exact degree sixteen.

\subsection{Final variable counts and exact projections}

In the parent 221-coordinate chart the removed private coordinates are
\[
 145,146,147,148,149,150,215,216,217,218,
\]
and the removed rows are
\[
 102,103,104,105,106,107,153,154,155,156.
\]
The checker verifies that none of these coordinates occurs in a
retained row. The one-certificate source has $221-10+1=212$ variables
and 151 rows; the split source has 213 variables and 152 rows.

The same five integral projections along directions in the common
kernel of the quadratic Hessians, followed by the same affine
consequence, remove six syntax coordinates. At every step the actual
pivot slope is $1$ or $-1$, the restoration has integer coefficients,
and the new pooling expressions are invariant along the projected
directions. The independent checker expands the shears directly rather
than using the producer's directional-derivative method. Hence the final
counts are 206 and 207.

Both equations share exactly 205 coordinates $x_0,\ldots,x_{204}$,
the same 144 quadratic residuals and the same sixteen request tuples.
For $P_{206}$ the only new coordinate is $z=x_{205}$; for $P_{207}$
they are $z_0=x_{205},z_1=x_{206}$. Appendix B specifies all those
residuals and request tuples literally, so its two short formulas define
both complete integer polynomials without importing any unspecified
universal polynomial.
